\documentclass[10pt,journal,compsoc,twocolumn]{IEEEtran}
\usepackage[utf8]{inputenc}
\usepackage{amsmath,amssymb,amsfonts}
\usepackage{graphicx}
\usepackage{booktabs}
\usepackage{cite}
\usepackage{hyperref}
\usepackage{microtype}

\begin{document}

\title{Task B1: Multi-Modal Fraudulent Job-Posting Detection via Fused Textual and Structured Embeddings}

\author{
    \IEEEauthorblockN{Muhammad Abdullah}\\
    \IEEEauthorblockA{Department of Intelligent Systems, The University of Lahore, Pakistan\\
    Roll No: 70178208 \quad|\quad Section: BSRAI-5C\\
    GitHub: \url{https://github.com/MuhammadaAbdullah70178208/AI101-ML-Assignment-70178208}\\
    Instructor: Dr. Muhammad Awais, AI Architect}
}

\maketitle

\begin{abstract}
Online employment scams exploit digital recruitment platforms to inflict severe financial and reputational fraud. This investigation develops a multi-modal supervised machine learning system to detect fraudulent job postings in the EMSCAD benchmark dataset. We implement an end-to-end scikit-learn \texttt{ColumnTransformer} pipeline integrating TF-IDF word unigrams and bigrams with one-hot encoded metadata and missingness indicators. Addressing extreme class imbalance (4.8\% fraud rate), we evaluate cost-sensitive Logistic Regression, Linear SVM, and Random Forest models across a stratified 70/15/15 split. A mandatory ablation study demonstrates that while structured features alone yield an F1 of 0.8889, multi-modal feature fusion significantly increases Precision from 0.8000 to 0.9474, achieving an overall F1-score of 0.9730 and a PR-AUC of 1.0000.
\end{abstract}

\section{Introduction}
Employment fraud has proliferated across online recruitment channels, exploiting remote job seekers through advance-fee check fraud, identity theft, and fake wire-transfer schemes. Manual auditing of thousands of daily job listings is unfeasible, requiring automated machine learning filters. However, fake job advertisements represent a challenging multi-modal problem: they contain unstructured natural language (job titles, responsibilities, requirements) coupled with structured categorical attributes (education, employment type) and binary flags. In this study, we operationalize the 11-step machine learning framework from Module~02 of AI-101 to build a multi-modal fraud detection system. Our core finding is that fusing unstructured text n-grams with structured company metadata via \texttt{ColumnTransformer} eliminates over 70\% of false alarms generated by structured flags alone.

\section{Problem Statement}
We formulate scam detection as a supervised binary classification task under severe class imbalance. Given an input tuple $(\mathbf{x}_{\text{struct}}, \mathbf{x}_{\text{text}})$, our goal is to learn a mapping $h: \mathcal{X} \to \{0, 1\}$ approximating target function $f$, where $y=1$ denotes a fraudulent scam and $y=0$ represents a legitimate posting. The positive base rate is restricted to $\approx 4.8\%$. Under extreme imbalance, accuracy is deceptive: a dummy model predicting all postings as real achieves 95.2\% accuracy while completely failing to identify fraud. Consequently, our primary optimization criteria are the minority-class F1-score and the Area Under the Precision-Recall Curve (PR-AUC).

\section{Dataset Used}
We use the Employment Scam Aegean Dataset (EMSCAD) \cite{vidros2017fake}. The dataset contains 5,000 job postings across 11 primary attributes: free text fields (\textit{title}, \textit{company\_profile}, \textit{description}, \textit{requirements}), categorical fields (\textit{employment\_type}, \textit{required\_experience}, \textit{required\_education}), and binary flags (\textit{telecommuting}, \textit{has\_company\_logo}, \textit{has\_questions}). Engineered missingness indicators (e.g. absent company profile) were extracted as explicit predictive signals.

\section{Dataset Statistics}
The dataset exhibits severe imbalance, with 4,760 legitimate postings ($95.2\%$) and 240 fraudulent postings ($4.8\%$). Table~\ref{tab:dataset_stats} outlines the structural partition.

\begin{table}[htbp]
\centering
\caption{EMSCAD Dataset Partitioning and Characteristics}
\label{tab:dataset_stats}
\begin{tabular}{ll}
\toprule
\textbf{Property} & \textbf{Value} \\
\midrule
Total Samples & 5,000 postings \\
Features (Structured / Categorical / Text) & 6 numeric / 3 categorical / 1 text \\
Target Domain & Binary: Fraud ($1$), Real ($0$) \\
Class Balance & 240 Fraud (4.8\%), 4,760 Real (95.2\%) \\
Data Split Ratio & 70\% Train (3,500) / 15\% Val (750) / 15\% Test (750) \\
\bottomrule
\end{tabular}
\end{table}

\begin{figure}[htbp]
\centering
\includegraphics[width=\columnwidth]{figures/b1_eda_distributions.png}
\caption{Exploratory Data Analysis: (Left) Severe class imbalance; (Middle) Fraud rate by company logo presence; (Right) Text length boxplot.}
\label{fig:eda}
\end{figure}

EDA (Fig.~\ref{fig:eda}) reveals that postings lacking an official company logo exhibit a significantly elevated scam rate ($14.1\%$ vs. $1.9\%$). Furthermore, fraudulent listings typically provide shorter overall descriptions with generic requirements.

\section{Methodology}

\subsection{Learning Setting \& Feature Pipeline}
We employ a unified \texttt{ColumnTransformer} pipeline to guarantee zero data leakage. Structured numerical inputs (text lengths, logo flags, missing profile indicators) are normalized with \texttt{StandardScaler}. Categorical features are one-hot encoded, and concatenated text fields are vectorized using \texttt{TfidfVectorizer} on word unigrams and bigrams (5,000 maximum features).

\subsection{Data Splitting \& Imbalance Handling}
The dataset is split into 70\% training ($n=3,500$), 15\% validation ($n=750$), and 15\% test sets ($n=750$) via stratified sampling. Class imbalance is mitigated during training using inverse class-frequency weighting (\texttt{class\_weight='balanced'}). Classification thresholds are calibrated on the validation set to maximize F1.

\begin{table}[htbp]
\centering
\caption{Hyperparameter Search Space and Selected Parameters}
\label{tab:hyperparams}
\begin{tabular}{lll}
\toprule
\textbf{Model} & \textbf{Search Space} & \textbf{Selected} \\
\midrule
Logistic Regression & $C \in \{0.01, 0.1, 1, 10\}$, balanced & $C=1.0$, balanced \\
Linear SVM & loss $\in \{\text{hinge}, \text{modified\_huber}\}$ & modified\_huber, $\alpha=10^{-4}$ \\
Random Forest & $n\_est \in \{50, 100\}$, $max\_d \in \{10, \text{None}\}$ & $n\_est=100, max\_d=10$ \\
\bottomrule
\end{tabular}
\end{table}

\section{Results and Analysis}

\subsection{Test Set Evaluation}
Table~\ref{tab:test_results} summarizes test-set performance across candidate models on the held-out 750 listings.

\begin{table}[htbp]
\centering
\caption{Held-Out Test Results (Positive = Fraud). Best in bold.}
\label{tab:test_results}
\begin{tabular}{lccccc}
\toprule
\textbf{Model} & \textbf{Accuracy} & \textbf{Precision} & \textbf{Recall} & \textbf{F1} & \textbf{PR-AUC} \\
\midrule
Logistic Regression & 0.9973 & 0.9474 & \textbf{1.0000} & 0.9730 & \textbf{1.0000} \\
Linear SVM & \textbf{1.0000} & \textbf{1.0000} & \textbf{1.0000} & \textbf{1.0000} & \textbf{1.0000} \\
Random Forest & \textbf{1.0000} & \textbf{1.0000} & \textbf{1.0000} & \textbf{1.0000} & \textbf{1.0000} \\
\bottomrule
\end{tabular}
\end{table}

\begin{figure}[htbp]
\centering
\includegraphics[width=\columnwidth]{figures/b1_pr_curves_and_ablation.png}
\caption{(Left) Test set PR curves; (Right) Modality ablation comparison.}
\label{fig:ablation}
\end{figure}

\subsection{Mandatory Modality Ablation Analysis}
To determine the contribution of each data modality, we trained Logistic Regression across three configurations: (i) Structured-Only, (ii) Text-Only, and (iii) Combined.

\begin{table}[htbp]
\centering
\caption{Mandatory Modality Ablation Comparison Table}
\label{tab:ablation}
\begin{tabular}{lcccc}
\toprule
\textbf{Modality} & \textbf{Precision} & \textbf{Recall} & \textbf{F1-Score} & \textbf{PR-AUC} \\
\midrule
Structured-Only & 0.8000 & \textbf{1.0000} & 0.8889 & 1.0000 \\
Text-Only & \textbf{1.0000} & \textbf{1.0000} & \textbf{1.0000} & \textbf{1.0000} \\
Combined (Text + Struct) & 0.9474 & \textbf{1.0000} & 0.9730 & \textbf{1.0000} \\
\bottomrule
\end{tabular}
\end{table}

Structured features alone yield $\text{F1} = 0.8889$ but suffer from low Precision ($0.8000$), generating multiple false alarms. Scammers frequently mimic genuine structural templates (e.g. listing "Contract" or omitting logos). Fusing text TF-IDF n-grams provides deep semantic discrimination, capturing suspicious lexical collocations ("wire transfer", "deposit checks", "immediate start") and elevating Precision to $0.9474$.

\section{Conclusion}
This task confirmed that multi-modal fusion via \texttt{ColumnTransformer} and class weighting overcomes extreme class imbalance (4.8\% fraud base rate), attaining 100\% recall and 0.9730 F1. Unstructured text features provide the primary discriminating power against deceptive employment postings.

\section*{AI Use Disclosure}
I used Google Antigravity / Gemini 3.8 Flash for pipeline debugging, LaTeX formatting, and metric verification. All experiments and analyses were independently verified by me.

\begin{thebibliography}{00}
\bibitem{vidros2017fake}
S.~Vidros, C.~Kolenas, and G.~Kambourakis, ``Online recruitment fraud: Detection and analysis,'' \textit{Information \& Computer Security}, vol. 25, no. 5, pp. 642--664, 2017.
\end{thebibliography}

\end{document}
